\section*{الفصل \ref{ch:g7:motion-average-speed} --- مخطّطات الحركة والسرعة المتوسّطة}

\begin{solution}{exo:g7:motion-average-speed:1}
$v = d/t$؛ والوحدتان \unit{km/h} و\unit{m/s}. المسافة:
$d = v \times t$؛ والزمن: $t = d/v$.
\end{solution}

\begin{solution}{exo:g7:motion-average-speed:2}
العبّارة: $45 \div 3 = \qty{15}{km/h}$. والأرنب: $100 \div 8 =
\qty{12.5}{m/s}$.
\end{solution}

\begin{solution}{exo:g7:motion-average-speed:3}
$5 \times 3.6 = \qty{18}{km/h}$؛ \ و$72 \div 3.6 = \qty{20}{m/s}$.
\end{solution}

\begin{solution}{exo:g7:motion-average-speed:4}
الخطّ المستقيم: \omterm{def:g6:describing-motion:uniform}{حركة منتظمة}؛ والأشدّ ميلًا: أسرع؛ والمستوي:
متوقّف. وأمّا الانحناء الذي يستوي تدريجيًّا فيدلّ على تباطؤ.
\end{solution}

\begin{solution}{exo:g7:motion-average-speed:5}
من الدقيقة $35$ إلى $60$ --- الجزء الثالث: فهو أقلّ الأجزاء
الصاعدة ميلًا، يُقرأ مباشرة من ميله الألين.
\end{solution}

\begin{solution}{exo:g7:motion-average-speed:6}
منتظمة --- خطّ مستقيم واحد. وتكسب \qty{2}{km} كل نصف ساعة:
$v = \qty{4}{km/h}$.
\end{solution}

\begin{solution}{exo:g7:motion-average-speed:7}
المسافة الكلّية على الزمن الكلّي، بما فيه التوقّفات. النزهة
كلها: $12 \div 4 = \qty{3}{km/h}$. والمشي وحده: $12 \div 3 =
\qty{4}{km/h}$.
\end{solution}

\begin{solution}{exo:g7:motion-average-speed:8}
على مسافتين متساويتين، يدوم الجزء البطيء أطول ببساطة --- فيُقضى
من ساعة الرحلة نصيب أكبر عيشًا بالوتيرة البطيئة، فتتكلّم
الوتيرة البطيئة بحصة أكبر من \omterm{def:g3:sound-around-us:sound}{صوت} الرحلة.
\end{solution}

\begin{solution}{exo:g7:motion-average-speed:9}
تخمين المنتصف: \qty{22.5}{km/h}. وبأمانة: الذهاب $30 \div 15 =
2$ ساعة، والعودة $30 \div 30 = 1$ ساعة؛ والمجموع \qty{60}{km}
في $3$ ساعات: \qty{20}{km/h}. فالنصف البطيء ملك ساعتين من
الثلاث.
\end{solution}

\begin{solution}{exo:g7:motion-average-speed:10}
المخطّط: مستقيم إلى (\qty{30}{min}، \qty{20}{km})؛ ثم استواء
حتى الدقيقة $45$؛ ثم مستقيم أشدّ ميلًا إلى (\qty{60}{min}،
\qty{35}{km}). والمتوسّط: \qty{35}{km} في ساعة واحدة ---
\qty{35}{km/h}.
\end{solution}

\begin{solution}{exo:g7:motion-average-speed:11}
$d = 340 \times 6 = \qty{2040}{m}$ --- نحو \qty{2}{km}.
والقاعدة القديمة: $6 \div 3 = 2$ كيلومتر --- فتتّفق معها، لأن
\qty{340}{m/s} تعني كيلومترًا كل ثلاث ثوانٍ تقريبًا.
\end{solution}

\begin{solution}{exo:g7:motion-average-speed:12}
النصف الأول: \qty{30}{km} \omterm{def:g7:motion-average-speed:formula}{بسرعة} \qty{30}{km/h} تكلّف ساعة $1$
بالضبط. ومتوسّط \qty{60}{km/h} على \qty{60}{km} يسمح بزمن كلّي
مقداره ساعة $1$ بالضبط --- وقد أُنفقت إلى آخر ثانية. فالمتبقّي
\qty{30}{km} يلزم أن يستغرق $0$ ساعة: فلا \omterm{def:g7:motion-average-speed:formula}{سرعة}، مهما بلغت من
البطولة، تصنع ذلك المتوسّط؛ فقد ذهبت الميزانية.
\end{solution}

\begin{solution}{pb:g7:motion-average-speed:1}
\textbf{1.} سير نشيط ($30$ دقيقة)، فتوقّف $15$ دقيقة (تسليم)،
فسير ألين ($30$ دقيقة)، فتوقّف ثانٍ مدّته $15$ دقيقة --- جولتان
وطلبان.
\textbf{2.} \qty{9}{km} في نصف ساعة: \qty{18}{km/h}.
\textbf{3.} $15 - 9 = \qty{6}{km}$ في نصف ساعة:
\qty{12}{km/h}.
\textbf{4.} \qty{15}{km} في $1.5$ ساعة: \qty{10}{km/h}.
\textbf{5.} $t = 12 \div 24 = 0.5$ ساعة: $30$ دقيقة.
\textbf{6.} $20$ دقيقة ثلث ساعة: $d = 18 \times \tfrac{1}{3} =
\qty{6}{km}$.
\textbf{7.} $25$ دقيقة هي $\tfrac{25}{60}$ ساعة؛ و$v = 10
\div \tfrac{25}{60} = \qty{24}{km/h}$.
\textbf{8.} المسافات: $12 + 6 + 10 = \qty{28}{km}$؛ والأزمنة:
$30 + 20 + 25 = 75$ دقيقة $= 1.25$ ساعة؛ ومتوسّط الدوران:
$28 \div 1.25 = \qty{22.4}{km/h}$.
\textbf{9.} $37 \div 3 \approx \qty{12.3}{km/h}$ --- أدنى من
$20$ بكثير: فلا مكافأة (وتوقّفاتها عمل شريف، لكن القاعدة
تعدّها).
\textbf{10.} ماكس: النصف السريع $12 \div 30 = 0.4$ ساعة؛
والنصف البطيء $12 \div 15 = 0.8$ ساعة؛ والمجموع \qty{24}{km} في
$1.2$ ساعة: \qty{20}{km/h} بالضبط --- لا \emph{فوق} $20$: فلا
مكافأة.
\textbf{11.} عدده $22.5$ حسب متوسّط \emph{العددين}، كأنّ كل
وتيرة تملك نصف الساعة. وفي الحقيقة ملك النصفُ البطيء $0.8$ من
ساعاته البالغة $1.2$ --- ثلثَي \omterm{def:g3:sound-around-us:sound}{صوت} الرحلة --- فجرّ المتوسّط
الحقيقي إلى $20$.
\textbf{12.} مثلًا: تُدفع المكافأة على \omterm{def:g7:motion-average-speed:average}{السرعة المتوسّطة}
\emph{أثناء الدوران} --- المسافة الكلّية مقسومة على زمن السير
وحده، باستثناء أجزاء المخطّط المستوية. فهي \omterm{def:g6:measurement-in-science:measuring}{تقيس} الوتيرة التي
تحفظها العجلات فعلًا، فتكافئ السير السريع لا التسليمات
المتخطّاة.
\textbf{13.} ``الجزء المستقيم وتيرة ثابتة؛ والجزء المستوي
توقّف والساعة ما زالت تدور؛ والجزء الأشدّ ميلًا وتيرة أسرع ---
فالميل هو \omterm{def:g7:motion-average-speed:formula}{السرعة}. ولا تحسب أبدًا متوسّط أعداد \omterm{def:g7:motion-average-speed:formula}{السرعة}: بل احسب
متوسّط الرحلة --- المسافة الكلّية على الزمن الكلّي --- وإلا
سخرت الساعات البطيئة من الورقة.''
\end{solution}
